% Part: normal-modal-logic
% Chapter: axioms-systems
% Section: normal-logics

\documentclass[../../../include/open-logic-section]{subfiles}

\begin{document}

\olfileid[id]{nml}{prf}{nor}

\olsection{Logika Modal Normal}

Tidak setiap himpunan !!{formula}s modal dapat dengan mudah dicirikan sebagai
himpunan !!{formula}s yang dapat diturunkan dari suatu himpunan aksioma. Kita
menginginkan logika modal yang berperilaku baik. Pertama-tama, semua yang dapat
kita !!{derive} dalam logika proposisional klasik tentu harus tetap
!!{derivable}, dengan memperhitungkan bahwa !!{formula}s kini juga dapat
memuat \iftag{prvBox}{$\Box$\iftag{prvDiamond}{ dan~}{}}{}
\iftag{prvDiamond}{$\Diamond$}{}. Untuk itu, kita mensyaratkan agar suatu
logika modal memuat semua instansiasi tautologis dan tertutup di bawah modus
ponens.

\begin{defn}
  Suatu \emph{logika modal} merupakan himpunan~$\Sigma$ dari !!{formula}s
  modal yang
  \begin{enumerate}
  \item memuat semua tautologi;
  \item tertutup di bawah substitusi, yakni jika $!A\in\Sigma$ dan
    $!D_1$, \dots, $!D_n$ merupakan !!{formula}s, maka
    \[
    \SSubst{!A}{\subst{!D_1}{p_1}, \dots, \subst{!D_n}{p_n}}\in\Sigma;
    \]
  \item tertutup di bawah \emph{modus ponens}, yakni jika $!A$ dan
    $!A\lif !B\in\Sigma$, maka $!B\in\Sigma$.
  \end{enumerate}
\end{defn}

Untuk menggunakan semantik relasional bagi logika modal, kita juga harus
mensyaratkan agar semua !!{formula}s yang valid dalam semua model modal
tercakup. Ternyata, syarat ini terpenuhi begitu semua instansiasi
\Ax{K}\iftag{prvDiamond}{ dan~\Dual{}}{} !!{derivable}, dan setiap kali
!!a{formula}~$!A$ !!{derivable}, demikian pula~$\Box !A$. Logika modal
yang memenuhi kondisi-kondisi ini disebut \emph{normal}. (Tentu saja ada pula
logika modal nonnormal, tetapi model relasional biasa tidak memadai baginya.)

\begin{defn}
  Suatu logika modal~$\Sigma$ \emph{normal} jika logika itu memuat
  \begin{align*}
    \tag{\Ax{K}} & \Box(p \lif q) \lif (\Box p \lif \Box q),
    \iftag{prvDiamond}{\\
      \tag{\Dual} & \Diamond p \liff \lnot\Box\lnot p}{}
  \end{align*}
  dan tertutup di bawah \emph{nesesitasi}, yakni jika $!A\in\Sigma$, maka
  $\Box !A\in\Sigma$.
\end{defn}

Perhatikan bahwa sementara implikasi tautologis cukup ``terperinci'' untuk
mempertahankan \emph{kebenaran pada suatu dunia}, aturan~\Nec{} hanya
mempertahankan \emph{kebenaran dalam suatu model} (dan karena itu juga
validitas pada suatu kerangka atau kelas kerangka).

\begin{prop}\ollabel{prop:rk}
  Setiap logika modal normal tertutup di bawah aturan~\RK,
  \begin{prooftree}
    \AxiomC{$!A_1 \lif (!A_2 \lif \cdots (!A_{n-1} \lif !A_n)\cdots)$}
    \RightLabel{\RK}
    \UnaryInfC{$\Box!A_1 \lif (\Box!A_2 \lif \cdots (\Box!A_{n-1}
      \lif \Box!A_n)\cdots).$}
  \end{prooftree}
\end{prop}

\begin{proof}
  Melalui induksi pada~$n$. Jika $n=1$, aturan tersebut tidak lain dari~\Nec,
  dan setiap logika modal normal tertutup di bawah~\Nec.

  Sekarang andaikan hasilnya berlaku bagi $n-1$; kita tunjukkan bahwa hasil itu
  berlaku bagi~$n$.

  Andaikan
  \begin{align*}
  & !A_1 \lif (!A_2 \lif \cdots (!A_{n-1} \lif !A_n)\cdots) \in \Sigma
  \intertext{Menurut hipotesis induksi, kita mempunyai}
  & \Box!A_1 \lif (\Box!A_2 \lif \cdots \Box(!A_{n-1} \lif !A_n)\cdots)
  \in \Sigma
  \intertext{Karena $\Sigma$ suatu logika modal normal, logika itu memuat
    semua instansiasi~$\Ax{K}$, khususnya}
  & \Box(!A_{n-1} \lif !A_n) \lif (\Box!A_{n-1} \lif \Box!A_n) \in \Sigma
  \intertext{Dengan menggunakan modus ponens dan instansiasi tautologis yang sesuai,
    kita memperoleh}
  & \Box!A_1 \lif (\Box!A_2 \lif \cdots (\Box!A_{n-1}
  \lif \Box!A_n)\cdots) \in \Sigma.
  \end{align*}
\end{proof}

\begin{prop}\ollabel{prop:notDiamondBot}
  Setiap logika modal normal~$\Sigma$ memuat~$\lnot\Diamond\lfalse$.
\end{prop}

\begin{prob}
  Buktikan \olref[nml][prf][nor]{prop:notDiamondBot}.
\end{prob}

\begin{prop}
  Misalkan $!A_1$, \dots, $!A_n$ merupakan !!{formula}s. Maka terdapat logika
  modal normal terkecil~$\Sigma$ yang memuat semua instansiasi dari
  $!A_1$, \dots, $!A_n$.
\end{prop}

\begin{proof}
  Bagi $!A_1$, \dots, $!A_n$ yang diberikan, definisikan~$\Sigma$ sebagai
  irisan semua logika modal normal yang memuat semua instansiasi dari
  $!A_1$, \dots, $!A_n$. Keluarga yang diiris tidak kosong, sebab $\Frm[L]$,
  himpunan semua !!{formula}s, merupakan logika modal semacam itu.
\end{proof}

\begin{defn}
Logika modal normal terkecil yang memuat $!A_1$, \dots, $!A_n$ disebut suatu
\emph{sistem modal} dan dilambangkan dengan $\Log{K} !A_1 \dots !A_n$.
Logika modal normal terkecil dilambangkan dengan~\Log{K}.
\end{defn}

\end{document}
